% One mock midterm in the announced format (75 minutes, short answers and problems that mix
% topics up to and including Grover's algorithm). Each of the four questions comes in two
% versions; solutions follow each version.

\clearpage
\noindent\textbf{Study copy.} Pages 1--53 above are the scribes exactly as submitted; the mock
exam below is not part of them. Do not bring this copy to the exam: only the submitted scribes are
allowed. The page numbers in the solutions are pages of the submitted scribes.
\vspace{0.2in}

\noindent\textbf{\large Mock Midterm}

\vspace{0.1in}

\noindent 75 minutes, 4 questions, 25 points each. Open scribes only. Each question has two
versions. For one practice exam, pick one version of each question (for example all Version 1,
then all Version 2, or mix them: $2^4 = 16$ different exams). Cover the solution, which follows
each version, until you have written your answer.
\vspace{0.1in}

\qslot{1}{Short answers}

\begin{question}[Short answers]
    Answer each part in two or three lines.
    \begin{enumerate}[label=(\alph*)]
        \item $X \geq 0$ and $\expectation{X} = 5$. Bound $\prob{X \geq 20}$. If also
              $\variance{X} = 4$, give a better bound.
        \item The birthday-paradox proof bounds $\prob{X = 0}$ with Chebyshev's inequality. Why not
              with the Chernoff bound?
        \item Is $\frac{1}{2}\left(\ket{00} + \ket{01} + \ket{10} - \ket{11}\right)$ entangled?
              Prove your answer.
        \item Why can no quantum circuit map $\ket{x_1 x_2}$ to $\ket{x_1 \oplus x_2}$? How does the
              oracle $U_f$ get around this?
        \item Deutsch--Jozsa uses $1$ query, and a deterministic algorithm needs $2^{n-1} + 1$. Is
              this an exponential separation from randomized algorithms? Is Simon's?
    \end{enumerate}
\end{question}
\begin{msolution}
    \uses{\scr{thm:markov}; \scr{thm:chebyshev}; \scr{lem:pairwise-0903}; \scr{thm:bell-0915};
          \scr{lem:xor-0915}; \scr{thm:dj-0924}; \scr{thm:simon-det-lb-1001}}
    (a) Markov: $\prob{X \geq 20} \leq \sfrac{5}{20} = \sfrac{1}{4}$. Chebyshev: $X \geq 20$
    implies $\abs{X - 5} \geq 15$, so $\prob{X \geq 20} \leq \sfrac{4}{15^2} = \sfrac{4}{225}
    \approx 0.018$.

    (b) Chernoff needs mutually independent indicators. The pair indicators $X_{ij}$ are only
    pairwise independent: if persons $1, 2$ and persons $2, 3$ share a birthday, so do $1, 3$.
    Pairwise independence is exactly what the variance needs, so Chebyshev applies.

    (c) Write the state as $a\ket{00} + b\ket{01} + c\ket{10} + d\ket{11}$. A product state has
    $ad = \alpha_0\alpha_1\beta_0\beta_1 = bc$. Here $ad = -\sfrac{1}{4}$ and $bc = \sfrac{1}{4}$,
    so the state is entangled.

    (d) Gates are unitary, so they are invertible and preserve norms. The map is not one-to-one:
    $\ket{00}$ and $\ket{11}$ would both go to $\ket{0}$, so $\ket{00} - \ket{11}$, of norm
    $\sqrt{2}$, would go to $0$. The oracle keeps the input:
    $U_f\ket{x}\ket{y} = \ket{x}\ket{y \oplus f(x)}$ permutes the basis states and is its own
    inverse, so it is unitary.

    (e) No. A randomized algorithm queries $11$ random inputs and answers ``balanced'' if two
    answers differ; it errs with probability at most $2^{-10}$, for every $n$. So the gap is
    $1$ against a constant. Simon's is exponential: every randomized algorithm with success
    probability $\sfrac{2}{3}$ needs more than $\sfrac{\sqrt{2^n}}{2}$ queries, against $\bigO{n}$
    quantum queries.
\end{msolution}

\begin{question}[Short answers]
    Answer each part in two or three lines.
    \begin{enumerate}[label=(\alph*)]
        \item A random walk on the path $1 - 2 - 3$ starts at vertex $1$. What is the stationary
              distribution? Does $p_t$ converge to it? Why?
        \item In the leader election algorithm, what goes wrong if each candidate contacts only
              $\ceil{\sqrt{n}}$ referees instead of $\ceil{\sqrt{n \ln n}}$?
        \item Which of $U = \frac{1}{\sqrt{2}}\begin{pmatrix} 1 & i \\ i & 1 \end{pmatrix}$ and
              $V = \begin{pmatrix} 1 & 1 \\ 0 & 1 \end{pmatrix}$ is unitary?
        \item A measurement-free quantum algorithm $\mathcal{A}$ succeeds with probability
              $\sfrac{1}{100}$. How many applications of $\mathcal{A}$ or $\mathcal{A}^{-1}$ does
              amplitude amplification need, against the expected number of classical repetitions?
              Why must $\mathcal{A}$ contain no measurement?
        \item In $\set{0,1}^n$ with $x \cdot y = \sum_i x_i y_i \bmod 2$, can a nonzero vector be
              orthogonal to itself? What does this mean for the measured $y$'s in Simon's
              algorithm?
    \end{enumerate}
\end{question}
\begin{msolution}
    \uses{\scr{thm:stationary-0910}; \scr{thm:convergence-0910}; \scr{lem:disjoint-0908};
          \scr{thm:leader-atmost-0908}; \scr{thm:unitary-0915}; \scr{thm:ampamp-1008};
          \scr{thm:orth-ind-0917}; \scr{thm:indep-1001}}
    (a) Degrees $1, 2, 1$ and $m = 2$, so $\pi = \left(\frac{1}{4}, \frac{1}{2}, \frac{1}{4}\right)$.
    The path is bipartite ($\set{1,3}$ and $\set{2}$), so it does not converge: $p_t$ alternates
    between $(0, 1, 0)$ and $\left(\frac{1}{2}, 0, \frac{1}{2}\right)$.

    (b) Two referee sets are disjoint with probability about $e^{-k^2/n} = e^{-1} \approx 0.37$
    instead of $\sfrac{1}{n}$, so two candidates often share no referee and both become leaders.
    The bound on $\prob{\text{more than one leader}}$ becomes $\sfrac{2\ln^2 n}{e} > 1$: useless.

    (c) $U^\dagger U = \frac{1}{2}\begin{pmatrix} 1 & -i \\ -i & 1 \end{pmatrix}\begin{pmatrix} 1 & i \\ i & 1 \end{pmatrix} = I$,
    so $U$ is unitary (remember to conjugate). $V$ is not: $V\ket{1} = \ket{0} + \ket{1}$ has norm
    $\sqrt{2}$.

    (d) $\sin\theta = \sqrt{\sfrac{1}{100}} = 0.1$, $\theta \approx 0.1002$,
    $k = \floor{\sfrac{\pi}{(4\theta)}} = 7$, success $\sin^2(15\theta) \approx 0.995$, cost
    $2k + 1 = 15$ applications, against $100$ expected classical repetitions. The operator
    $Q = -\mathcal{A}Z_0\mathcal{A}^{-1}Z_f$ needs $\mathcal{A}^{-1}$, and a measurement cannot be
    undone.

    (e) Yes: $11 \cdot 11 = 1 + 1 = 0$. So ``orthogonal implies independent'' (proved over
    $\mathbb{R}$ and $\mathbb{C}$ using $\norm{v}^2 > 0$) fails over $\set{0,1}$. In Simon's
    algorithm all $y$'s satisfy $s \cdot y = 0$, and $s$ itself may be among them; their
    independence comes from the probability bound of the Oct 1 exercise, not from orthogonality.
\end{msolution}

\qslot{2}{Quantum query algorithms}

\begin{question}[One circuit, two problems]
    Let $f \colon \set{0,1}^3 \to \set{0,1}$, $f(x) = x_1 \oplus x_3$.
    \begin{enumerate}[label=(\alph*)]
        \item Run the circuit $H^{\otimes 3}$, $U_f$, $H^{\otimes 3}$ on $\ket{000}\ket{-}$: write
              the state after the query and after the last Hadamards. What is measured?
        \item What does the Deutsch--Jozsa algorithm conclude about $f$? What does the
              Bernstein--Vazirani algorithm conclude? Are both conclusions correct?
        \item Prove that every deterministic classical algorithm that finds $s$ from queries to
              $f(x) = s \cdot x$ needs $n$ queries.
    \end{enumerate}
\end{question}
\begin{msolution}
    \uses{\scr{thm:dj-0924}; \scr{thm:bv-0929}; \scr{lem:signsum-0929}}
    (a) After the first Hadamards and the query (phase kickback), the first register is
    \[\frac{1}{\sqrt{8}} \sum_{x} (-1)^{x_1 \oplus x_3} \ket{x}
        = \frac{1}{\sqrt{8}} \left(\ket{000} - \ket{001} + \ket{010} - \ket{011} - \ket{100} + \ket{101} - \ket{110} + \ket{111}\right).\]
    With $s = 101$, $f(x) = s \cdot x$. After the last Hadamards the amplitude of $\ket{y}$ is
    $\frac{1}{8} \sum_x (-1)^{(s \oplus y) \cdot x}$, which is $1$ for $y = s$ and $0$ otherwise
    (sum of signs). So $101$ is measured with probability $1$.

    (b) The outcome is not $000$, so Deutsch--Jozsa says ``balanced''. Bernstein--Vazirani says
    $s = 101$. Both are correct: $f$ is $1$ on exactly the $4$ inputs with $x_1 \neq x_3$.

    (c) Fix a deterministic algorithm. Each query $x$ returns one bit $s \cdot x$, a linear
    equation in the unknown $s$. The adversary answers $0$ to every query. After $q < n$ queries,
    the strings $s$ with $s \cdot x = 0$ for every queried $x$ form a subspace of $\set{0,1}^n$
    of dimension at least $n - q \geq 1$, so it contains $0^n$ and some $s \neq 0^n$. Both agree
    with every answer, and the algorithm gives the same output on both, so it is wrong on one.
    Querying $10\dots0, \dots, 0\dots01$ shows that $n$ queries suffice.
\end{msolution}

\begin{question}[Simon's algorithm]
    The function $f \colon \set{0,1}^3 \to \set{0,1}^3$ is $2$-to-$1$ with an unknown period $s$:
    \begin{center}
        \begin{tabular}{c|cccccccc}
            $x$    & 000 & 001 & 010 & 011 & 100 & 101 & 110 & 111 \\ \hline
            $f(x)$ & 101 & 010 & 000 & 110 & 000 & 110 & 101 & 010
        \end{tabular}
    \end{center}
    \begin{enumerate}[label=(\alph*)]
        \item In one run of Simon's circuit the second register is measured as $010$. Give the state
              of the first register, then after $H^{\otimes 3}$. Which $y$ can be measured, and with
              what probability?
        \item Two runs give $y = 111$ and $y = 001$. Find $s$, and say how to check it.
        \item With what probability are the $y$'s of two runs linearly independent? Compare with
              the scribe's bound.
        \item Why does every classical algorithm need $\bigOm{\sqrt{2^n}}$ queries? Two or three
              sentences.
    \end{enumerate}
\end{question}
\begin{msolution}
    \uses{\scr{thm:simon-run-0929}; \scr{thm:simon-0929}; \scr{thm:indep-1001};
          \scr{lem:collision-1001}; \scr{thm:simon-rand-lb-1001}}
    (a) The inputs with $f(x) = 010$ are $001$ and $111$, so the first register is
    $\frac{1}{\sqrt{2}}\left(\ket{001} + \ket{111}\right)$. After $H^{\otimes 3}$ the amplitude of
    $\ket{y}$ is $\frac{1}{4}(-1)^{001 \cdot y}\left(1 + (-1)^{110 \cdot y}\right)$: it is
    $\pm\frac{1}{2}$ when $110 \cdot y = 0$ and $0$ otherwise. So $y \in \set{000, 001, 110, 111}$,
    each with probability $\sfrac{1}{4}$.

    (b) $s \cdot 111 = 0$ and $s \cdot 001 = 0$ give $s_3 = 0$ and $s_1 = s_2$, so $s \in \set{000,
    110}$. Check with two classical queries: $f(000) = 101 = f(110)$, so $s = 110$.

    (c) The $y$'s are uniform in $V = \set{y \suchthat s \cdot y = 0}$, which has $4$ elements. The
    first must be nonzero ($\sfrac{3}{4}$) and the second outside $\set{000, y_1}$ ($\sfrac{2}{4}$):
    probability $\sfrac{3}{8}$. The exercise gives $\left(1 - \frac{1}{2}\right)\left(1 - \frac{1}{4}\right)
    = \sfrac{3}{8} \geq \sfrac{1}{4}$, the same.

    (d) A classical algorithm learns about $s$ only from a collision, two queries $x$ and
    $x \oplus s$. With $k$ queries there are $\binom{k}{2}$ pairs, each a collision with
    probability about $\sfrac{2}{2^n}$ for a random $s$, so the collision probability is at most
    $\sfrac{k(k-1)}{2^n}$: small unless $k = \bigOm{\sqrt{2^n}}$ (birthday paradox). By Yao's
    principle this holds for randomized algorithms too.
\end{msolution}

\qslot{3}{Probability and randomized algorithms}

\begin{question}[Random sampling]
    Random sampling makes $m$ independent uniformly random queries into an array of $n$ slots.
    \begin{enumerate}[label=(\alph*)]
        \item Find the expected number of slots that are never accessed.
        \item Let $m = 2n \ln n$. Show that some slot is never accessed with probability at most
              $\sfrac{1}{n}$.
        \item Let $m = 3n \ln n$. Show that with probability at least $1 - n^{-11}$ no slot is
              accessed $18 \ln n$ times or more.
    \end{enumerate}
\end{question}
\begin{msolution}
    \uses{\scr{thm:linearity-0827}; \scr{thm:markov}; \scr{thm:whp}; \scr{thm:chernoff-0901}}
    (a) Let $Y_i = 1$ if slot $i$ is never accessed. A query misses slot $i$ with probability
    $1 - \sfrac{1}{n}$, so $\expectation{Y_i} = \left(1 - \sfrac{1}{n}\right)^m$, and by linearity
    the expected number of empty slots is
    $n\left(1 - \sfrac{1}{n}\right)^m \leq n e^{-m/n}$, using $1 - x \leq e^{-x}$.

    (b) For $m = 2n \ln n$ this is at most $n e^{-2 \ln n} = \sfrac{1}{n}$. The number $Y$ of empty
    slots is $\geq 0$, so by Markov $\prob{Y \geq 1} \leq \expectation{Y} \leq \sfrac{1}{n}$.
    (Equivalently: union bound over the $n$ slots.)

    (c) Fix a slot; the number of accesses $Z$ is a sum of $3n \ln n$ independent indicators with
    probability $\sfrac{1}{n}$, so $\mu = 3 \ln n$. Since $18 \ln n = 6\mu$, the third Chernoff form
    gives $\prob{Z \geq 18 \ln n} \leq 2^{-18 \ln n} = n^{-18 \ln 2} \leq n^{-12}$. Union bound
    over $n$ slots: at most $n^{-11}$.
\end{msolution}

\begin{question}[A random walk]
    Vertices $1, 2, 3$ form a triangle and vertex $4$ is joined to vertex $3$.
    \begin{enumerate}[label=(\alph*)]
        \item Write the transition matrix $P$.
        \item Find the stationary distribution $\pi$ and verify $\pi P = \pi$.
        \item Does the walk converge to $\pi$ from every start? Why?
        \item Starting at vertex $4$, compute $p_1$ and $p_2$.
    \end{enumerate}
\end{question}
\begin{msolution}
    \uses{\scr{lem:evolution-0910}; \scr{thm:stationary-0910}; \scr{thm:convergence-0910}}
    (a) Degrees $2, 2, 3, 1$:
    \[P = \begin{pmatrix}
            0            & \tfrac{1}{2} & \tfrac{1}{2} & 0            \\
            \tfrac{1}{2} & 0            & \tfrac{1}{2} & 0            \\
            \tfrac{1}{3} & \tfrac{1}{3} & 0            & \tfrac{1}{3} \\
            0            & 0            & 1            & 0
        \end{pmatrix}.\]

    (b) $m = 4$, so $\pi = \left(\frac{2}{8}, \frac{2}{8}, \frac{3}{8}, \frac{1}{8}\right)$. Check,
    vertex by vertex: $\frac{1}{8} + \frac{1}{8} = \frac{1}{4}$;
    $\frac{1}{8} + \frac{1}{8} = \frac{1}{4}$; $\frac{1}{8} + \frac{1}{8} + \frac{1}{8} = \frac{3}{8}$;
    $\frac{3}{8} \cdot \frac{1}{3} = \frac{1}{8}$.

    (c) Yes: the graph is connected, and the triangle is an odd cycle, so it is not bipartite.

    (d) $p_0 = (0,0,0,1)$, $p_1 = (0, 0, 1, 0)$, $p_2 = \left(\frac{1}{3}, \frac{1}{3}, 0, \frac{1}{3}\right)$.
\end{msolution}

\qslot{4}{Grover and lower bounds}

\begin{question}[Grover with numbers]
    Let $N = 8$ with exactly one solution.
    \begin{enumerate}[label=(\alph*)]
        \item Track the eight amplitudes through one Grover iteration ($Z_f$, then inversion about
              the mean). What is the success probability after one iteration?
        \item Find $\theta$, the number of iterations $k$ of the scribe's exercise, and the success
              probability for $k = 2$ and $k = 3$. Explain why $k = 3$ is worse.
        \item Now the number of solutions $a$ is unknown, $1 \leq a \leq 4$. Give an algorithm with
              $\bigO{\sqrt{N}}$ queries that succeeds with probability at least $\sfrac{3}{8}$,
              and make the success probability at least $0.99$.
    \end{enumerate}
\end{question}
\begin{msolution}
    \uses{\scr{lem:diffusion-1006}; \scr{cor:hk-1008}; \scr{thm:iterations-1008}; \scr{thm:randomk-1008}}
    (a) Start with all amplitudes $\sfrac{1}{\sqrt{8}}$. $Z_f$ makes the solution's amplitude
    $-\sfrac{1}{\sqrt{8}}$. The mean is $\frac{1}{8} \cdot \frac{7 - 1}{\sqrt{8}} = \frac{3}{4\sqrt{8}}$.
    Inversion about the mean sends $\alpha$ to $2 \cdot \text{mean} - \alpha$: the non-solutions
    become $\frac{1}{2\sqrt{8}}$ and the solution $\frac{5}{2\sqrt{8}}$. Success probability
    $\frac{25}{32} \approx 0.78$. Check: $7 \cdot \frac{1}{32} + \frac{25}{32} = 1$.

    (b) $\sin\theta = \sfrac{1}{\sqrt{8}}$, $\theta \approx 0.3614$ rad $\approx 20.7^\circ$.
    $\frac{\pi}{4\theta} - \frac{1}{2} \approx 1.67$, so $k = 2$. After $k$ iterations the success
    probability is $\sin^2\left((2k+1)\theta\right)$: for $k = 2$ the angle is $103.5^\circ$ and
    the probability is $\approx 0.945$ (at least $1 - \sfrac{1}{8}$, as the exercise
    guarantees); for $k = 3$ the angle is $144.9^\circ$ and the probability drops to
    $\approx 0.33$. Each iteration rotates by $2\theta$, and the third one rotates past
    $90^\circ$, away from the solution.

    (c) Choose $k$ uniformly from $\set{1, \dots, \floor{2\sqrt{8}}} = \set{1, \dots, 5}$, make
    $k$ iterations and measure. Since $a \leq \sfrac{N}{2}$, this succeeds with probability at
    least $\sfrac{3}{8}$, using at most $5 = \bigO{\sqrt{N}}$ queries. Check the result with one
    query; if it is not a solution, repeat with a new random $k$. After $r$ rounds the failure
    probability is at most $\left(\sfrac{5}{8}\right)^r \leq 0.01$ for $r = 10$.
\end{msolution}

\begin{question}[Search lower bounds]
    An array of $n$ slots is searched for $x$ by queries ``is $A[i] = x$?''.
    \begin{enumerate}[label=(\alph*)]
        \item Without a promise that $x$ is present, show that a deterministic algorithm that
              decides whether $x$ is present needs $n$ queries.
        \item With the promise, show that a randomized algorithm that finds $x$ with probability
              at least $\sfrac{2}{3}$ on every input needs at least $\sfrac{2n}{3} - 1$ queries.
        \item How many queries does Grover's algorithm make for $n = 2^{20}$, and why does (b) not
              contradict this?
    \end{enumerate}
\end{question}
\begin{msolution}
    \uses{\scr{thm:detsearch-0825}; \scr{lem:yao-1001}}
    (a) The adversary answers ``no'' to every query. If the algorithm stops after at most $n - 1$
    queries, some slot $j$ is unqueried. The array with $x$ in slot $j$ and the array with no $x$
    agree with all answers, and the algorithm gives the same output on both: wrong on one.

    (b) Let $\mu$ put $x$ in a uniformly random slot. A deterministic algorithm with $T$ queries
    queries a fixed sequence $q_1, \dots, q_T$ while the answers are ``no''. It finds $x$ with
    probability $\sfrac{T}{n}$; otherwise it outputs one fixed slot, right with probability at
    most $\sfrac{1}{n}$. So it succeeds with probability at most $\sfrac{(T+1)}{n}$ under $\mu$.
    If $T < \sfrac{2n}{3} - 1$, every such algorithm errs with probability more than $\sfrac{1}{3}$
    under $\mu$, and by Yao's principle every randomized algorithm with $T$ queries errs with
    probability more than $\sfrac{1}{3}$ on some input.

    (c) $\theta \approx \sfrac{1}{1024}$, $k = \floor{\sfrac{\pi}{(4\theta)}} \approx
    \floor{\frac{\pi}{4} \cdot 1024} = 804$ iterations, one query each. (b) is a bound for
    classical algorithms; quantum queries can be made in superposition, which is what Grover
    exploits. The quadratic speedup is the best possible: quantum search needs
    $\bigOm{\sqrt{N}}$ queries.
\end{msolution}
